\documentclass[oneside,numbers,spanish]{ezthesis}
\usepackage[utf8]{inputenc}
\usepackage{paralist}
\usepackage{listings}
\usepackage{url}

\lstset{language=C++,texcl=true}

%% # Opciones disponibles para el documento #
%%
%% Las opciones con un (*) son las opciones predeterminadas.
%%
%% Modo de compilar:
%%   draft            - borrador con marcas de fecha y sin im'agenes
%%   draftmarks       - borrador con marcas de fecha y con im'agenes
%%   final (*)        - version final de la tesis
%%
%% Tama'no de papel:
%%   letterpaper (*)  - tama'no carta (Am'erica)
%%   a4paper          - tama'no A4    (Europa)
%%
%% Formato de impresi'on:
%%   oneside          - hojas impresas por un solo lado
%%   twoside (*)      - hijas impresas por ambos lados
%%
%% Tama'no de letra:
%%   10pt, 11pt, o 12pt (*)
%%
%% Espaciado entre renglones:
%%   singlespace      - espacio sencillo
%%   onehalfspace (*) - espacio de 1.5
%%   doublespace      - a doble espacio
%%
%% Formato de las referencias bibliogr'aficas:
%%   numbers          - numeradas, p.e. [1]
%%   authoryear (*)   - por autor y a'no, p.e. (Newton, 1997)
%%
%% Opciones adicionales:
%%   spanish         - tesis escrita en espa'nol
%%
%% Desactivar opciones especiales:
%%   nobibtoc   - no incluir la bibiolgraf'ia en el 'Indice general
%%   nofancyhdr - no incluir "fancyhdr" para producir los encabezados
%%   nocolors   - no incluir "xcolor" para producir ligas con colores
%%   nographicx - no incluir "graphicx" para insertar gr'aficos
%%   nonatbib   - no incluir "natbib" para administrar la bibliograf'ia

%% Paquetes adicionales requeridos se pueden agregar tambi'en aqu'i.
%% Por ejemplo:
%\usepackage{subfig}
%\usepackage{multirow}

%% # Datos del documento #
%% Nota que los acentos se deben escribir: \'a, \'e, \'i, etc.
%% La letra n con tilde es: \~n.

\author{Roberto Castellor Morant}
\title{Sistema básico de almacenamiento en la nube usando Java RMI}
\degree{Grado de Ingeniería Informática}
\supervisor{Nombre de mi Asesor}
\institution{Universidad Nacional de Educación a Distancia}
\faculty{Grado de Ingeniería Informática}
\department{Departamento de Sistemas de Comunicación y Control}

%% # M'argenes del documento #
%% 
%% Quitar el comentario en la siguiente linea para austar los m'argenes del
%% documento. Leer la documentaci'on de "geometry" para m'as informaci'on.

\geometry{top=40mm,bottom=33mm,inner=40mm,outer=25mm}

%% El siguiente comando agrega ligas activas en el documento para las
%% referencias cruzadas y citas bibliogr'aficas. Tiene que ser *la 'ultima*
%% instrucci'on antes de \begin{document}.
\hyperlinking
\begin{document}

%% En esta secci'on se describe la estructura del documento de la tesis.
%% Consulta los reglamentos de tu universidad para determinar el orden
%% y la cantidad de secciones que debes de incluir.

%% # Portada de la tesis #
%% Mirar el archivo "titlepage.tex" para los detalles.
%% ## Construye tu propia portada ##
%% 
%% Una portada se conforma por una secuencia de "Blocks" que incluyen
%% piezas individuales de informaci'on. Un "Block" puede incluir, por
%% ejemplo, el t'itulo del documento, una im'agen (logotipo de la universidad),
%% el nombre del autor, nombre del supervisor, u cualquier otra pieza de
%% informaci'on.
%%
%% Cada "Block" aparece centrado horizontalmente en la p'agina y,
%% verticalmente, todos los "Blocks" se distruyen de manera uniforme 
%% a lo largo de p'agina.
%%
%% Nota tambi'en que, dentro de un mismo "Block" se pueden cortar
%% lineas usando el comando \\
%%
%% El tama'no del texto dentro de un "Block" se puede modificar usando uno de
%% los comandos:
%%   \small      \LARGE
%%   \large      \huge
%%   \Large      \Huge
%%
%% Y el tipo de letra se puede modificar usando:
%%   \bfseries - negritas
%%   \itshape  - it'alicas
%%   \scshape  - small caps
%%   \slshape  - slanted
%%   \sffamily - sans serif
%%
%% Para producir plantillas generales, la informaci'on que ha sido inclu'ida
%% en el archivo principal "tesis.tex" se puede accesar aqu'i usando:
%%   \insertauthor
%%   \inserttitle
%%   \insertsupervisor
%%   \insertinstitution
%%   \insertdegree
%%   \insertfaculty
%%   \insertdepartment
%%   \insertsubmitdate

\begin{titlepage}
  \TitleBlock{\scshape\insertinstitution}
  \TitleBlock[\bigskip]{\scshape\insertfaculty}
  \TitleBlock{\Huge\scshape\inserttitle}
  \TitleBlock{\scshape
    Memoria presentada por \insertauthor \\
    para la asignatura Sistemas Distribuidos de \insertdegree}
  \TitleBlock{DNI: 16041663A \\
  Correo electrónico: roberto.castellor@gmail.com}
  \TitleBlock{\insertsubmitdate}
  \TitleBlock[\bigskip]{\insertdepartment}
  \TitleBlock{\includegraphics[height=2cm]{uned.jpg}}
\end{titlepage}

%% Nota 1:
%% Se puede agregar un escudo o logotipo en un "Block" como:
%% y teniendo un archivo "escudo_uni.pdf", "escudo_uni.png" o "escudo_uni.jpg"
%% en alg'un lugar donde LaTeX lo pueda encontrar.

%% Nota 2:
%% Normalmente, el espacio entre "Blocks" se extiende de modo que el
%% contenido se reparte uniformemente sobre toda la p'agina. Este
%% comportamiento se puede modificar para mantener fijo, por ejemplo, el
%% espacio entre un par de "Blocks". Escribiendo:
%%   \TitleBlock{Bloque 1}
%%   \TitleBlock[\bigskip]{Bloque2}
%% se deja un espacio "grande" y de tama~no fijo entre el bloque 1 y 2.
%% Adem'as de \bigskip est'an tambi'en \smallskip y \medskip. Si necesitas
%% aun m'as control puedes usar tambi'en, por ejemplo, \vspace*{2cm}.




%% # Prefacios #
%% Por cada prefacio (p.e. agradecimientos, resumen, etc.) crear
%% un nuevo archivo e incluirlo aqu'i.
%% Para m'as detalles y un ejemplo mirar el archivo "gracias.tex".
%%\include{gracias}

%% # 'Indices y listas de contenido #
%% Quitar los comentarios en las lineas siguientes para obtener listas de
%% figuras y cuadros/tablas.
\tableofcontents
%\listoffigures
%\listoftables

%% # Cap'itulos #
%% Por cada cap'itulo hay que crear un nuevo archivo e incluirlo aqu'i.
%% Mirar el archivo "intro.tex" para un ejemplo y recomendaciones para
%% escribir.

\chapter{Introducción}

Este documento contiene la memoria de la práctica obligatoria de laboratorio correspondiente a la asignatura de \emph{Sistemas Distribuidos} del Grado de Ingeniería Informática. El propósito de la practica es el desarrollo de un software que implemente un sistema de almacenamiento de ficheros en la nube usando Java RMI.

El aplicativo consiste en tres ejecutables diferentes, el servidor levanta tres servicios, dos de ellos serán usados por los distintos usuarios del servicio distribuido mientras que el tercer servicio sera para uso interno de los servicios propios del servidor.

El repositorio publica dos objetos distribuidos para uso de los clientes o del servidor, estos objetos ofrecen los métodos necesarios para almacenar los ficheros en su sistema local de ficheros o enviarlos al sistema local de ficheros de los clientes.

El cliente sera la parte utilizada por el usuario para gestionar sus ficheros, publicara a su vez un objeto distribuido que se utilizara para almacenar ficheros en el disco local del cliente.

La interfaz de la aplicación se implementa mediante un menú en modo texto, los resultados de las operaciones se volcaran sobre la salida estándar. De cara a una depuración más eficiente del sistema se generará un fichero de log llamado ``log.txt"\ en la carpeta donde se ejecuten los procesos donde se podrá encontrar información adicional para depurar la aplicación.


\chapter{Casos de uso}

En este capitulo repasamos los principales casos de uso de la aplicación, caben analizar muchos más escenarios y en sucesivas iteraciones de desarrollo se podría ampliar el numero de casos de uso asi como logar un mayor nivel de detalle de los expuestos en el documento.

\section{Inicio del servidor}

\begin{compactenum}
	\item El caso de uso comienza cuando el usuario lanza el proceso servidor.
	\item El sistema levanta el registro para almacenar los objetos distribuidos
	\item El sistema crea el objeto para el servicio de datos
	\item El sistema publica en el registro RMI el servicio de datos
	\item El sistema crea el objeto para el servicio de autenticación
	\item El sistema publica en el registro RMI el servicio de autenticación
	\item El sistema crea el objeto para el servicio gestor
	\item El sistema publica en el registro RMI el servicio gestor
	\item El caso de uso esta completo, el sistema entra en modo escucha de comandos
\end{compactenum}
\includegraphics[width=0.8\linewidth]{inicio_servidor.png}

\section{Registrar/autenticar repositorio}

\begin{compactenum}
	\item El caso de uso comienza cuando el usuario del repositorio selecciona la opción de registrar repositorio
	\item El sistema solicita al usuario los datos de autenticación
	\item El sistema busca el servicio de autenticación en el registro RMI
	\item El sistema comprueba los parámetros de entrada con el objeto remoto
	\item El sistema publica el servicio operador cliente
	\item El sistema publica el servicio operador servidor
	\item El sistema entra en modo consulta para operaciones con el repositorio activo
\end{compactenum}
\includegraphics[width=0.8\linewidth]{registro_rep.png}

\section{Registrar/autenticar cliente}

\begin{compactenum}
	\item El caso de uso comienza cuando el usuario del cliente selecciona la opción de registrar cliente
	\item El sistema solicita al usuario los datos de autenticación
	\item El sistema busca el servicio de autenticación en el registro RMI
	\item El sistema comprueba los parámetros de entrada con el objeto remoto
	\item El sistema publica los objetos remotos para operar con el servidor en el registro
	\item El sistema entra en modo consulta para operaciones con el cliente activo
\end{compactenum}
\includegraphics[width=0.8\linewidth]{registro_cli.png}

\section{Subir fichero a la nube}
\begin{compactenum}
	\item El caso de uso comienza cuando el usuario selecciona la opción subir fichero
	\item El sistema pregunta al usuario cual es el fichero que quiere subir a la nube
	\item El usuario selecciona el nombre del fichero
	\item El sistema crea un objeto serializable en memoria con el fichero a enviar
	\item El sistema busca el objeto distribuido gestor
	\item El sistema consulta al servicio gestor por la dirección del servicio operador cliente del repositorio
	\item El sistema gestor consulta al servicio de datos por el id del repositorio
	\item El sistema busca el objeto distribuido cliente operador del repositorio
	\item El sistema envia el objeto de tipo Fichero Serializable al repositorio mediante el servicio operador cliente
	\item El servicio operador-cliente guarda el fichero en disco con el metodo escribirEn del objeto Fichero
	\item El sistema notifica al sistema gestor el resultado del envío
	\item El servicio gestor actualiza el servicio de datos con el nuevo fichero
	\item El sistema notifica al usuario el resultado de la operación
\end{compactenum}
\includegraphics[width=\linewidth]{ds_subirfichero.png}

\section{Descargar fichero desde la nube}
\begin{compactenum}
\item El caso de uso comienza cuando el usuario selecciona la opcion de descargar fichero
\item El sistema pregunta al servicio gestor por los ficheros a los que tiene acceso el cliente
\item El sistema pregunta al usuario cual de los ficheros quiere descargar
\item El usuario selecciona un fichero de la lista
\item El sistema busca en el registro el objeto operador cliente del repositorio
\item El sistema pide al objeto operador cliente que envie el fichero mediante el servicio disco cliente
\item El servicio operador busca el servicio disco cliente
\item El servicio operador cliente envia el fichero a traves del metodo del servicio disco cliente
\item El servicio disco cliente guarda el fichero en el disco local del cliente
\end{compactenum}
\includegraphics[width=\linewidth]{ds_bajarfichero.png}

\section{Borrar fichero de la nube}
\begin{compactenum}
\item El caso de uso comienza cuando el usuario seleccion la opción de borrar fichero
\item El sistema pregunta al servicio gestor por los ficheros del usuario
\item El sistema pregunta al usuario cual es el fichero que quiere borrar
\item El sistema pide al servicio gestor la url del objeto distribuido del servicio operador-cliente de su repositorio
\item El sistema busca en el registro RMI el servicio operador-cliente del repositorio donde esta el fichero
\item El sistema pide al servicio operador-cliente que borre el fichero
\item El sistema notifica al servicio gestor del resultado de la operación de borrado
\item El servicio gestor notifica al servicio de datos que el fichero se ha borrado fisicamente
\end{compactenum}


\section{Compartir fichero con otro cliente}

\chapter{Consideraciones de diseño}

\section{Estructura de directorios}

Las Carpeta comun contiene las clases e interfaces comunes al resto del proyecto.

\begin{compactitem}
	\item \textbf{Cliente}, clase con la información de los clientes, el usuario, el id y listas con los ficheros del cliente
	\item \textbf{Fichero}, clase proporcionada por el equipo docente para enviar los ficheros
	\item \textbf{GUI}, clase con métodos genéricos para mostrar por pantalla información y recoger las entradas del usuario
	\item \textbf{Logger}, clase con métodos para volcar a un fichero log de la aplicación y de esta forma no mezclarse con la interfaz del usuario
	\item \textbf{Metadato}, clase con información referente a los ficheros, tendrá el repositorio, el id del cliente propietario del fichero, etc.
	\item \textbf{Repositorio}, clase con datos de los repositorios registrados en el sistema, los id, el usuario y las url de los objetos remotos que publica el repositorio.
	\item \textit{ServicioAutenticacionInterface}, interfaz del servicio de autenticación.
	\item \textit{ServicioDatosInterface}, interfaz del servicio de datos
	\item \textit{ServicioGestorInterface}, interfaz del servicio gestor
	\item \textit{ServicioClOperadorInterface}, interfaz del servicio cliente-operador
	\item \textit{ServicioSrOperadorInterface}, interfaz del servicio servidor-operador
	\item \textit{ServicioDiscoClienteInterface}, interfaz del servicio de disco del cliente
	\item \textbf{Utils}, clase con métodos genéricos para el uso por parte de la aplicación
\end{compactitem}

La carpeta servidor, contiene el código necesario para lanzar el servidor y la implementación del servicio de datos, el servicio de autenticación y el servicio gestor.
\begin{compactitem}
	\item \textbf{Servidor}, clase que lanza el servidor, crea los objetos remotos, los publica en el registro RMI y genera el menu para operar con el servidor
	\item \textbf{ServicioAutenticacionImpl}, clase que implementa el servicio de autenticación
	\item \textbf{ServicioDatosImpl}, clase que implementa el servicio de datos
	\item \textbf{ServicioGestorImpl}, clase que implementa el servicio gestor
\end{compactitem}

La carpeta repositorio contiene el código necesario para lanzar los repositorios y la implementación de los servicios operador-cliente y operador-servidor.

\begin{compactitem}
	\item \textbf{Repositorio}, clase que lanza el repositorio y genera el menú para operar con el y publica los objetos remotos.
	\item \textbf{ServicioClOperadorImpl}, clase que implementa el servicio cliente-operador
	\item \textbf{ServicioSrOperadorImpl}, clase que implementa el servicio servidor-operador
\end{compactitem}

La carpeta cliente contiene el código necesario para lanzar los clientes y la implementación de la interfaz remota del servicio de disco del cliente

\begin{compactitem}
	\item \textbf{Cliente}, clase que lanza el cliente, genera el menú para poder operar con el y publica los objetos remotos
	\item \textbf{ServicioDiscoClienteImpl}, clase que implementa el servicio de disco del cliente a través del cual se recibirán ficheros de la nube
\end{compactitem}

\section{Interfaces remotas}

\includegraphics[width=0.75\linewidth]{interfaces.png}

\section{Objetos modelo de datos}
\includegraphics[width=0.75\linewidth]{definicion_datos.png}

El modelo de datos esta basado en la api Collections de Java.

Cada elemento del sistema sera identificado por un id numérico que formara parte del objeto como atributo.

En el servicio de datos se almacenara la siguiente información:
\begin{compactitem}
\item Lista con todos los clientes registrados en la aplicación.
\item Relación de nombre de usuario con id de clientes en una HashMap$<$String, Integer$>$
\item Relación de id de usuario con el objeto de tipo Cliente que tiene toda la información referente a cada cliente.
\item Variable sesioncli con el ultimo identificador de cliente asignado, a cada cliente nuevo se le asignara como identificador el valor de esta variable y esta se incrementara en espera del próximo cliente.
\item Lista con todos los repositorios registrados en la aplicación.
\item Relación de nombre de usuario con id de repositorio en una HashMap$<$String, Integer$>$
\item Relación de id de usuario con el objeto de tipo Repositorio que tiene toda la información referente a cada cliente.
\item Variable sesionrep con el ultimo identificador de repositorio asignado, a cada repositorio nuevo se le asignara como identificador el valor de esta variable y esta se incrementara en espera del próximo cliente.
\item Variable numerica que se utilizara para decidir a que repositorio asignar cada nuevo cliente
\end{compactitem}

La clase Cliente del paquete comun tiene la información propia del cliente y dos mapas indexadas por el nombre del fichero, una con los metadatos de los ficheros cuyo propietario es el cliente y otra con los metadatos de los ficheros que se han compartido con el cliente. Las operaciones con los datos de los ficheros se realizaran mediante llamadas a metodos de esta clase y en ningún momento accediendo directamente a las collection que albergan los datos, (ver apartado \ref{sec:consideraciones_datos})

Esta información se persistirá a disco cuando se cierre el sistema para guardar el estado del servidor.

\section{Entorno del servidor}

El entorno del servidor usara las siguientes clases

\centerline{\includegraphics[width=0.5\linewidth]{servidor.png}}


\chapter{Ejemplos de uso}

En este capitulo mostraremos ejemplos de uso de la aplicación, para los ejemplos he registrado en el servidor cuatro repositorios y siete clientes, cada cliente tendra asociados varios ficheros.
En el siguiente diagrama se pueden apreciar los elementos de la sesión mostrada.

\includegraphics[width=0.8\linewidth]{escenario_pruebas.png}

\section{Operaciones del servidor}
\begin{lstlisting}[language=bash,frame=single,texcl=true,basicstyle=\small]
============Menu Servidor============
1-Listar Clientes
2-Listar Repositorios
3-Listar Parejas Repositorio-Cliente
4-Salir
============Menu Servidor============
Seleccione una opcion: 1   
-- Clientes registrados en el sistema --
Cliente: 0-cliente1
Cliente: 1-cliente2
Cliente: 2-cliente3
Cliente: 3-cliente4
Cliente: 4-cliente5
Cliente: 5-cliente6
Cliente: 6-cliente7
============Menu Servidor============
1-Listar Clientes
2-Listar Repositorios
3-Listar Parejas Repositorio-Cliente
4-Salir
============Menu Servidor============
Seleccione una opcion: 2
-- Repositorios registrados en el sistema --
Repositorio: 0-repositorio1
Repositorio: 1-repositorio2
Repositorio: 2-repositorio3
Repositorio: 3-repositorio4
============Menu Servidor============
1-Listar Clientes
2-Listar Repositorios
3-Listar Parejas Repositorio-Cliente
4-Salir
============Menu Servidor============
Seleccione una opcion: 3
-- Clientes - repositorio registrados en el sistema --
Cliente: 0-cliente1 -> 0-repositorio1
Cliente: 1-cliente2 -> 1-repositorio2
Cliente: 2-cliente3 -> 2-repositorio3
Cliente: 3-cliente4 -> 3-repositorio4
Cliente: 4-cliente5 -> 0-repositorio1
Cliente: 5-cliente6 -> 1-repositorio2
Cliente: 6-cliente7 -> 2-repositorio3
============Menu Servidor============
1-Listar Clientes
2-Listar Repositorios
3-Listar Parejas Repositorio-Cliente
4-Salir
============Menu Servidor============
Seleccione una opcion: 
\end{lstlisting}


\section{Operaciones del repositorio}
Este es un ejemplo de las operaciones disponibles en los repositorios.
\subsection{Registro nuevo repositorio}
\begin{lstlisting}[language=bash,frame=single,texcl=true,basicstyle=\small]
============Menu Repositorio============
1-Registrar nuevo repositorio
2-Autenticarse en el sistema
3-Salir
============Menu Repositorio============
Seleccione una opcion: 1
Introduzca el usuario: repositorio1
Introduzca la contraseña: rep1
Introduzca la contraseña de nuevo: rep1
============ repositorio0 ============
1-Listar Clientes
2-Listar ficheros del Cliente
3-Salir
============ repositorio0 ============
Seleccione una opcion: 
\end{lstlisting}
\subsection{Operaciones repositorio}
\begin{lstlisting}[language=bash,frame=single,texcl=true,basicstyle=\small]
============ repositorio0 ============
1-Listar Clientes
2-Listar ficheros del Cliente
3-Salir
============ repositorio0 ============
Seleccione una opcion: 1   
Listado de clientes en el repositorio
Cliente: 0-cliente1
Cliente: 4-cliente5
============ repositorio0 ============
1-Listar Clientes
2-Listar ficheros del Cliente
3-Salir
============ repositorio0 ============
Seleccione una opcion: 2
Clientes validos: 
0 4 
Seleccione un cliente: 0
Fichero: test1.txt
Fichero: test2.txt
============ repositorio0 ============
1-Listar Clientes
2-Listar ficheros del Cliente
3-Salir
============ repositorio0 ============
Seleccione una opcion: 
\end{lstlisting}

\section{Cliente}
Ejemplo de uso del cliente
\subsection{Registro nuevo cliente}
\begin{lstlisting}[language=bash,frame=single,texcl=true,basicstyle=\small]
============Menu Cliente============
1-Registrar nuevo cliente
2-Autenticarse en el sistema
3-Salir
============Menu Cliente============
Seleccione una opcion: 1
Introduzca usuario: cliente3
Introduzca contraseña: cli3
Introduzca la contraseña de nuevo: cli3
============-cliente3-============
1-Subir fichero
2-Bajar fichero
3-Borrar fichero
4-Compartir ficheros
5-Listar ficheros
6-Listar clientes del sistema
7-Salir
============-cliente3-============
Seleccione una opcion: 
\end{lstlisting}
\subsection{Operación subir fichero}
\begin{lstlisting}[language=bash,frame=single,texcl=true,basicstyle=\small]
============-cliente3-============
1-Subir fichero
2-Bajar fichero
3-Borrar fichero
4-Compartir ficheros
5-Listar ficheros
6-Listar clientes del sistema
7-Salir
============-cliente3-============
Seleccione una opcion: 1
Introduzca fichero a enviar: test1.txt
Fichero enviado correctamente a la nube
\end{lstlisting}
\subsection{Operación descargar fichero}
\begin{lstlisting}[language=bash,frame=single,texcl=true,basicstyle=\small]
============-cliente3-============
1-Subir fichero
2-Bajar fichero
3-Borrar fichero
4-Compartir ficheros
5-Listar ficheros
6-Listar clientes del sistema
7-Salir
============-cliente3-============
Seleccione una opcion: 2
test1.txt 
Seleccione el fichero que quiere recibir: test1.txt
Escribiendo fichero en test1.txt
\end{lstlisting}
\subsection{Operación borrar un fichero}
\begin{lstlisting}[language=bash,frame=single,texcl=true,basicstyle=\small]
============-cliente3-============
1-Subir fichero
2-Bajar fichero
3-Borrar fichero
4-Compartir ficheros
5-Listar ficheros
6-Listar clientes del sistema
7-Salir
============-cliente3-============
Seleccione una opcion: 3
test1.txt 
Seleccione el fichero que quiere eliminar: test1.txt
\end{lstlisting}
\subsection{Otras operaciones}
Listado de ficheros del cliente y de los clientes del sistema.

\begin{lstlisting}[language=bash,frame=single,texcl=true,basicstyle=\small]
============-cliente1-============
1-Subir fichero
2-Bajar fichero
3-Borrar fichero
4-Compartir ficheros
5-Listar ficheros
6-Listar clientes del sistema
7-Salir
============-cliente1-============
Seleccione una opcion: 5   
Fichero: test1.txt
Fichero: test2.txt
============-cliente1-============
1-Subir fichero
2-Bajar fichero
3-Borrar fichero
4-Compartir ficheros
5-Listar ficheros
6-Listar clientes del sistema
7-Salir
============-cliente1-============
Seleccione una opcion: 6
Cliente: 0-cliente1
Cliente: 1-cliente2
Cliente: 2-cliente3
Cliente: 3-cliente4
Cliente: 4-cliente5
Cliente: 5-cliente6
Cliente: 6-cliente7
\end{lstlisting}

%%\include{conclu}

\appendix
%% Cap'itulos incluidos despues del comando \appendix aparecen como ap'endices
%% de la tesis.
\chapter{Herramientas}

Para realizar esta practica se han utilizado estas herramientas:
\begin{compactitem}
	\item OpenJDK Java 1.8 http://openjdk.java.net/
	\item Ant para construir los ejecutables del proyecto http://ant.apache.org/
	\item Eclipse para edición del código http://www.eclipse.org/
	\item git como sistema para mantener el control de la evolución del código y permitir desarrollar desde varios PC https://git-scm.com/
	\item draw.io, herramienta en la nube para generar graficos https://www.draw.io/
	\item Latex, entorno tex para realizar la memoria
\end{compactitem}

\chapter{Riesgos, consideraciones y conclusiones}

El proceso para desarrollo de la practica ha sido un proceso incremental por lo que algunas decisiones de diseño pueden no ser todo lo correctas, en el código fuente entregado esta también el histórico y la evolución del código, el histórico se puede consultar usando la herramienta de control de versiones \textbf{git}.

\section{Revisión de decisiones de diseño}

\begin{enumerate}\label{sec:consideraciones_datos}
	\item La clase cliente del paquete común necesita una refactorización, en muchos métodos remotos estamos enviando toda la información del cliente, eso no esta optimizado y en sucesiones iteraciones del desarrollo habría que optimizarlo para enviar solo la información necesaria
	\item En la aplicación no se comprueba si los servicios están publicados correctamente y se da por supuesto que todos los objetos remotos están operativos. En siguientes iteraciones del proceso de desarrollo habría que ver que implicaciones tendría tener un nivel mas fino de control de errores.
	\item El repositorio crea una carpeta con su id para almacenar los ficheros, simplemente es para facilitar las pruebas con varios repositorios lanzados desde la misma carpeta.
\end{enumerate}

%\include{apendiceB}
%\include{apendiceC}

%% Incluir la bibliograf'ia. Mirar el archivo "biblio.bib" para m'as detales
%% y un ejemplo.
\bibliography{biblio}

\end{document}
